\subsection{Supervised learning method}
A summary of the results of the Random Forest model is listed in Table \ref{tab:ml_results_small}. An $R^2$ of 0.96 is found on the train data for both the full
and the reduced model, and an $R^2$ of 0.81 for the full model on the test data, and 0.82 for the reduced model. The Mean Absolute Percentage Error for the
reduced model is slightly lower for the reduced model, on the test data. This implies the full model is slightly over-fitting on the training data.
\begin{longtable}{lccc}
\caption{Performance metrics for the Random forest models using all features and reduced features.} \label{tab:ml_results_small} \\
\hline
Model & $R^2$ & MAE (m) & MAPE (\%) \\
\hline
\endfirsthead
\hline
Model & $R^2$ & MAE (m) & MAPE (\%) \\
\hline
\endhead
Full (train) & 0.96 & 831.94 & 5.62\\
Full (test) & 0.81 & 1888.69 & 12.48\\
Reduced (train) & 0.96 & 831.89 & 5.62\\
Reduced (test) & 0.82 & 1859.49 & 12.16\\
\hline
\end{longtable}


Figure \ref{tab:feature-importance} displays the most important features for predicting TSP path length, according to the Random Forest model that is used.
The average path length between all nodes in the graph of the road network, $n$ (the number of locations in the TSP) and the area of the neighborhood the TSP is
in are by far the most important features.

It is obvious why the number of locations is an important predictor.
The average path length, the radius, as well as the area say something about how spread out an area is.
All three of these features are at the top of the most important features.
The average edge length also carries some importance.
This value might mostly imply something about how straight or curvy the roads are.
A road with a lot of curves has many short edges, and a long straight road consists of a few long edges.
The mean, variance and maximum value of the node degree also provide some predictive value.
It is difficult to interpret the importance of the features, since it might be that a feature that carries some importance does not directly affect TSP path length,
but is correlated with some other variable that is not known which in turn has an effect on TSP path length.
\begin{figure}[H]
	\centering
	\caption{The 15 most important features in the reduced model.}
	\includegraphics[width=0.9\textwidth]{../project/plots/feature_importances_reduced.png}
	\label{tab:feature-importance}
\end{figure}

Table \ref{tab:ml_results} shows per-neighborhood Mean Absolute Percentage errors, for the test data, comparing the full and the reduced model.
The model is able to predict TSP path length with below 5\% for some areas, but for other areas the prediction errors are close to or above 30\%.
Again, similarly to the results for the BHH formula, the predictions for the Stad van de Zon quarter have the highest MAPE, of all areas.
It seems that the Random Forest model is also not able to provide accurate results for areas with strong clustering of locations, as was the BHH formula,
even though some features were included in the Random Forest model that could provide some information about this clustering.
A potential reason for this might be that this clustering causes larger variation in the TSP path length, especially for smaller numbers of visited locations,
which would imply that any approximation would yield higher prediction errors in such area.

Overall, the Random Forest model provides estimates with MAPE almost half of the MAPE of the predictions of the BHH formula (with estimated value for $\beta$).
This is a strong improvement in performance, requiring little data collection to extend this to completely new areas.
